%%==========================================================================
\section{Heuristic Algorithms}
\label{sec:heuristics}
%%==========================================================================

We study four heuristic families in the present report.  The greedy+insertion
construction provides the main baseline, while the repair-style refinement
and the bounded local-search refinement are used to improve the same initial
solution before the final schedule is accepted.  The Dijkstra-based space-time
heuristic and the ACO variant are retained as alternative search strategies,
but the experimental results below focus on the deterministic constructive
methods because they remain the most stable under full maintenance pressure.

\subsection{Timeline simulator}

The core primitive is a function \textsc{GetTimeline}$(j, [f_1,\ldots,f_k])$
that evaluates the feasibility and cost of assigning an ordered flight sequence
to aircraft~$j$.  Maintenance triggers follow priority $D \succ C \succ B
\succ A$ per the counter logic of \Cref{sec:hierarchy}.  The simulator returns
\textsc{Infeasible} if any of the following occur:
\begin{itemize}
  \item A required ferry (repositioning) cannot complete before the next flight
    departure.
  \item A triggered maintenance check cannot complete before the next flight
    departure.
  \item Maintenance station capacity is exhausted at the current airport.
  \item The current time exceeds a flight's departure time.
\end{itemize}
Otherwise it returns the complete event timeline and total cost.

\subsection{Heuristic 1: Greedy + Insertion}

\paragraph{Phase 1 — Greedy.}
Flights are sorted by departure time.  Each flight $f$ is appended to the
route of the aircraft minimising incremental cost:
\begin{equation}
  j^* = \arg\min_{j \in \Pc}\;
  \mathrm{cost}\!\left(\textsc{GetTimeline}(j, R_j \Vert [f])\right).
\end{equation}
Infeasible assignments are skipped.

\paragraph{Phase 2 — Insertion repair.}
Up to 5 repair passes attempt to insert each unassigned flight into all
possible positions in all aircraft routes, accepting the least-cost feasible
insertion.  This refinement is deliberately conservative: it preserves the
greedy seed whenever no improving insertion exists, and it therefore acts as
a robust fallback that does not destabilise the constructive solution.

\paragraph{Refinement variant: bounded local search.}
In addition to the repair sweep, the implementation used for the benchmark
runs includes a bounded local-search refinement.  The method starts from the
greedy seed and then performs a limited number of relocation passes, moving a
flight from one aircraft route to another whenever the move improves the
current assignment while preserving maintenance feasibility.  A final
insertion sweep is then applied to remove any remaining gaps.  This variant is
more aggressive than the repair pass, but it is still controlled by the same
timeline simulator and therefore remains compatible with the same feasibility
logic as the baseline heuristic.

\paragraph{Complexity.}
$O(n_f^3 n_p)$ in the worst case (5 passes, $O(n_f)$ positions, $O(n_f)$
simulator cost).  In practice, significantly cheaper due to early termination
and sparse routes.

\subsection{Heuristic 2: Dijkstra-Based Space-Time Heuristic}

\paragraph{Network structure.}
The space-time network $G = (V, E)$ has nodes $(k, t)$ for airport $k$ at
time $t$ and arcs of five types: flight arcs (TrArcs), ground-wait arcs
(WArcs), repositioning arcs (DArcs), coupling/decoupling arcs (CDArcs), and
maintenance arcs.

\paragraph{Proxy costs.}
Arc costs are ranked to encode priority: flight arcs receive the lowest cost
(rank of actual cost among unserved flight arcs); CDArcs receive the second
lowest; ground arcs are proportional to wait duration; depot arcs are very high;
maintenance arcs are highest (to be traversed only when triggered).

\paragraph{Modified Dijkstra.}
The update rule is \emph{longest-distance first}: paths are updated when the
new route covers more flight arcs, or the same number at lower cost.  Arcs
served by the required number of aircraft are removed (cost~$\to\infty$).
The algorithm terminates at maintenance events or the horizon end.

\paragraph{Complexity.}
$O(n_p n_d n_f^2)$ using a binary heap.

{\sloppy
\subsection{Heuristic 3: Ant Colony Optimization (ACO)}

Each ant simulates an aircraft route.  The next node is selected according to:
\begin{equation}
  p_{ij} = \frac{\tau_{ij}\,(\eta_{ij})^\beta}
  {\sum_{u \in N(i)} \tau_{iu}\,(\eta_{iu})^\beta},
\end{equation}
where $\tau_{ij}$ is pheromone, $\eta_{ij} = 1/w(i,j)$ is heuristic
desirability, and $\beta \ge 1$ balances exploration vs.\ exploitation.
After each route is built, pheromone evaporates according to
\begin{equation}
  \tau_{ij} \leftarrow (1-\alpha)\tau_{ij} + \alpha\tau_0.
\end{equation}
Fully served arcs receive $\tau_{ij} = 0$ and this zero propagates backward
to all arcs that lead exclusively to the exhausted arc.
}

\paragraph{Why Dijkstra outperforms ACO.}
The space-time network has a strict temporal ordering (no back-edges in time).
This acyclicity removes the feedback loops that ACO exploits most effectively.
Dijkstra's farthest-distance criterion is near-optimal for acyclic networks,
which explains its consistent quality advantage over ACO on the tested
instances.  The empirical summary in Table~\ref{tab:heuristic_comparison}
therefore supports the interpretation that ACO is more exploratory but less
stable as a production solver for this particular temporal assignment
structure.

\subsection{Quality summary}

\begin{table}[t]
\centering\small
\caption{Empirical heuristic summary from the current benchmark suite
  (\texttt{ABCD\_cap\_bottleneck}, heavy density-$0.5$ cases,
  and dense density-$1.0$ cases; 8 instances total).}
\label{tab:heuristic_comparison}
\begin{tabular}{lccc}
\toprule
Property & Greedy+Ins. & Dijkstra & ACO \\
\midrule
Solver required? & No  & No  & No \\
Deterministic?   & Yes & Yes & No (seeded) \\
Avg. unassigned flights & 68.38 & 130.75 & 233.00 \\
Complete assignments (of 8) & 3 & 3 & 0 \\
Avg. wall time (s) & 3.800 & 11.384 & 121.024 \\
\bottomrule
\end{tabular}
\end{table}
